\documentclass[11pt,a4paper]{article}

% ============================================================
% Cours de mathématiques -- Première spécialité
% Chapitre : Suites numériques
% ============================================================

\usepackage[utf8]{inputenc}
\usepackage[T1]{fontenc}
\usepackage[french]{babel}
\usepackage{lmodern}
\usepackage{amsmath,amssymb,amsthm}
\usepackage{geometry}
\usepackage{fancyhdr}
\usepackage{tikz}
\usepackage[most]{tcolorbox}
\usepackage{enumitem}
\usepackage{array,tabularx}
\usepackage{multicol}
\usepackage{hyperref}
\usepackage{titlesec}
\usepackage{lastpage}

\geometry{left=2cm,right=2cm,top=2.1cm,bottom=2.2cm}
\setlength{\parindent}{0pt}
\setlength{\parskip}{0.45em}
\renewcommand{\arraystretch}{1.25}

\pagestyle{fancy}
\fancyhf{}
\lhead{Première générale -- Spécialité mathématiques}
\rhead{Suites numériques}
\cfoot{\thepage/\pageref{LastPage}}

\definecolor{bleuprepa}{RGB}{20,55,110}
\definecolor{vertprepa}{RGB}{20,105,70}
\definecolor{rougeprepa}{RGB}{160,35,35}
\definecolor{orangeprepa}{RGB}{180,100,20}
\definecolor{violetprepa}{RGB}{90,55,140}
\definecolor{grisclair}{RGB}{245,247,250}

\hypersetup{colorlinks=true,linkcolor=bleuprepa,urlcolor=bleuprepa}

\titleformat{\section}{\Large\bfseries\color{bleuprepa}}{}{0em}{\thesection.\;}
\titleformat{\subsection}{\large\bfseries\color{vertprepa}}{}{0em}{\thesubsection.\;}
\titleformat{\subsubsection}{\bfseries\color{orangeprepa}}{}{0em}{\thesubsubsection.\;}

\tcbset{
  enhanced,
  breakable,
  sharp corners=downhill,
  boxrule=0.7pt,
  colback=white,
  fonttitle=\bfseries,
  before skip=0.7em,
  after skip=0.7em
}

\newtcbtheorem[number within=section]{definitionbox}{Définition}{colframe=bleuprepa,colback=blue!3!white}{def}
\newtcbtheorem[number within=section]{propositionbox}{Propriété}{colframe=vertprepa,colback=green!3!white}{prop}
\newtcbtheorem[number within=section]{theorembox}{Théorème}{colframe=violetprepa,colback=violet!3!white}{thm}
\newtcbtheorem[number within=section]{examplebox}{Exemple corrigé}{colframe=orangeprepa,colback=orange!3!white}{ex}
\newtcbtheorem[number within=section]{exercisebox}{Exercice}{colframe=black!65,colback=grisclair}{exo}

\newtcolorbox{methodbox}[1][]{colframe=bleuprepa,colback=blue!2!white,title={Méthode #1}}
\newtcolorbox{rememberbox}[1][]{colframe=vertprepa,colback=green!3!white,title={À retenir #1}}
\newtcolorbox{warningbox}[1][]{colframe=rougeprepa,colback=red!3!white,title={Erreur classique #1}}
\newtcolorbox{proofbox}[1][]{colframe=black!65,colback=white,title={Démonstration #1}}
\newtcolorbox{correctionbox}[1][]{colframe=violetprepa,colback=violet!2!white,title={Correction #1}}

\newcommand{\N}{\mathbb{N}}
\newcommand{\R}{\mathbb{R}}
\newcommand{\suite}[1]{\left(#1_n\right)_{n\in\N}}
\newcommand{\abs}[1]{\left|#1\right|}

\begin{document}

\begin{titlepage}
\centering
\vspace*{1.5cm}
{\Huge\bfseries Cours complet de mathématiques\\[0.2cm]}
{\LARGE Première générale -- Spécialité mathématiques\\[0.8cm]}
{\Huge\color{bleuprepa}\bfseries Algèbre : suites numériques\\[0.4cm]}
\rule{0.8\textwidth}{0.5pt}\\[0.6cm]
{\Large Notion de suite, suites explicites et récurrentes, suites arithmétiques et géométriques, sens de variation, sommes géométriques, modélisations discrètes.}\\[1cm]
\begin{tcolorbox}[colframe=bleuprepa,colback=blue!2!white,width=0.9\textwidth,title=Objectifs du chapitre]
À la fin de ce chapitre, l'élève doit savoir :
\begin{itemize}
  \item comprendre qu'une suite est une fonction définie sur les entiers naturels ;
  \item calculer des termes d'une suite définie explicitement, par récurrence ou par un algorithme ;
  \item reconnaître, utiliser et démontrer les formules relatives aux suites arithmétiques et géométriques ;
  \item étudier le sens de variation d'une suite ;
  \item calculer une somme finie de termes consécutifs d'une suite géométrique ;
  \item modéliser une évolution discrète : population, placement financier, répétition d'un phénomène.
\end{itemize}
\end{tcolorbox}
\vfill
{\large Document professeur -- version longue, progressive et rigoureuse}\par
\end{titlepage}

\tableofcontents
\newpage

\section*{Préambule professeur}
\addcontentsline{toc}{section}{Préambule professeur}
Ce cours est rédigé pour une classe de Première générale suivant l'enseignement de spécialité mathématiques. Il adopte une exigence de rédaction inspirée de la classe préparatoire, mais reste conforme à l'esprit et aux méthodes attendues au lycée : calcul exact, raisonnement clair, modélisation, usage ponctuel de l'algorithmique et progressivité.

Le chapitre des suites est central, car il introduit une première manière de penser les phénomènes qui évoluent étape par étape. Contrairement aux fonctions de variable réelle, où la variable peut varier continûment, une suite décrit une grandeur observée à des instants entiers : année 0, année 1, année 2, etc. C'est le langage naturel des évolutions discrètes.

\begin{rememberbox}
Dans tout le chapitre, il faut insister sur trois idées :
\begin{enumerate}
  \item une suite est une liste ordonnée de nombres, indexée par un entier ;
  \item une formule explicite donne directement $u_n$ en fonction de $n$ ;
  \item une relation de récurrence donne un terme à partir du ou des précédents.
\end{enumerate}
\end{rememberbox}

\section{Notion générale de suite numérique}

\subsection{Définition d'une suite}

\begin{definitionbox}{Suite numérique}{suite}
Une \textbf{suite numérique} est une fonction définie sur une partie de l'ensemble des entiers naturels et à valeurs réelles.

Autrement dit, à chaque entier $n$ appartenant à un ensemble d'indices, on associe un nombre réel noté $u_n$.

On note généralement une suite :
\[
(u_n)_{n\in\N}, \qquad \text{ou plus simplement} \qquad (u_n).
\]
Le nombre $u_n$ s'appelle le \textbf{terme de rang $n$} de la suite.
\end{definitionbox}

\begin{examplebox}{Premiers exemples de suites}{premiers-exemples}
\begin{enumerate}
  \item La suite définie par $u_n=2n+1$ donne :
  \[
  u_0=1,\quad u_1=3,\quad u_2=5,\quad u_3=7.
  \]
  C'est la suite des nombres impairs positifs si l'on commence au rang $0$.

  \item La suite définie par $v_n=n^2$ donne :
  \[
  v_0=0,\quad v_1=1,\quad v_2=4,\quad v_3=9,\quad v_4=16.
  \]

  \item La suite définie par $w_0=5$ et $w_{n+1}=w_n+3$ donne :
  \[
  w_0=5,\quad w_1=8,\quad w_2=11,\quad w_3=14.
  \]
\end{enumerate}
\end{examplebox}

\begin{warningbox}
Il ne faut pas confondre $u_n$ et $(u_n)$.
\begin{itemize}
  \item $u_n$ désigne un nombre : le terme de rang $n$.
  \item $(u_n)$ désigne toute la suite, c'est-à-dire l'objet mathématique complet.
\end{itemize}
Dire « la suite $u_n$ » est fréquent à l'oral, mais dans une rédaction propre, on écrit plutôt « la suite $(u_n)$ ».
\end{warningbox}

\subsection{Rang initial : commencer à $0$ ou à $1$}

Une suite peut commencer au rang $0$, au rang $1$, ou parfois à un autre rang. Cela dépend du contexte.

\begin{examplebox}{Rang initial dans une situation concrète}{rang-initial}
Une population compte $10\,000$ individus en 2026. On note $P_n$ la population $n$ années après 2026.

Alors :
\[
P_0=10\,000.
\]
Le rang $0$ correspond à l'année de départ. Le rang $1$ correspond à 2027, le rang $2$ à 2028, etc.

Si l'on note au contraire $Q_n$ la population pendant l'année $n$, alors $Q_{2026}=10\,000$.
\end{examplebox}

\begin{rememberbox}
Avant de résoudre un problème de suites, il faut toujours identifier clairement :
\begin{itemize}
  \item ce que représente $n$ ;
  \item ce que représente $u_n$ ;
  \item le premier rang utilisé ;
  \item la valeur initiale.
\end{itemize}
\end{rememberbox}

\subsection{Suite définie par une formule explicite}

\begin{definitionbox}{Définition explicite}{explicite}
Une suite $(u_n)$ est définie par une \textbf{formule explicite} lorsque l'on peut calculer directement $u_n$ en fonction de $n$.

La forme générale est :
\[
 u_n=f(n),
\]
où $f$ est une expression dépendant de $n$.
\end{definitionbox}

\begin{examplebox}{Calcul de termes avec une formule explicite}{calcul-explicite}
Soit la suite $(u_n)$ définie pour tout $n\in\N$ par
\[
 u_n=3n^2-2n+1.
\]
Calculons les quatre premiers termes.
\[
 u_0=3\times 0^2-2\times 0+1=1,
\]
\[
 u_1=3\times 1^2-2\times 1+1=2,
\]
\[
 u_2=3\times 2^2-2\times 2+1=12-4+1=9,
\]
\[
 u_3=3\times 3^2-2\times 3+1=27-6+1=22.
\]
Donc les premiers termes sont :
\[
1,\;2,\;9,\;22.
\]
\end{examplebox}

\begin{methodbox}[Calculer un terme d'une suite explicite]
Pour calculer $u_k$ lorsque $u_n$ est donné explicitement :
\begin{enumerate}
  \item repérer la formule de $u_n$ ;
  \item remplacer chaque $n$ par $k$ ;
  \item effectuer les calculs en respectant les priorités opératoires ;
  \item conclure par une phrase.
\end{enumerate}
\end{methodbox}

\subsection{Suite définie par récurrence}

\begin{definitionbox}{Définition par récurrence}{recurrence}
Une suite est définie par \textbf{récurrence} lorsque l'on donne :
\begin{itemize}
  \item un ou plusieurs premiers termes ;
  \item une relation permettant de calculer un terme à partir du terme précédent, ou des termes précédents.
\end{itemize}

La forme la plus courante en Première est :
\[
\begin{cases}
 u_0=a,\\
 u_{n+1}=F(u_n).
\end{cases}
\]
\end{definitionbox}

\begin{examplebox}{Calcul de termes d'une suite récurrente}{calcul-recurrent}
Soit $(u_n)$ définie par
\[
\begin{cases}
 u_0=2,\\
 u_{n+1}=3u_n-1.
\end{cases}
\]
Calculons les premiers termes.
\[
 u_1=3u_0-1=3\times 2-1=5,
\]
\[
 u_2=3u_1-1=3\times 5-1=14,
\]
\[
 u_3=3u_2-1=3\times 14-1=41.
\]
Donc :
\[
 u_0=2,\quad u_1=5,\quad u_2=14,\quad u_3=41.
\]
\end{examplebox}

\begin{warningbox}
Dans une suite définie par récurrence, on ne peut généralement pas calculer directement $u_{50}$ sans calculer les termes précédents, sauf si l'on a d'abord trouvé une formule explicite.
\end{warningbox}

\subsection{Suite définie par un algorithme}

Une suite peut aussi être définie par un programme. En Première, on utilise souvent Python pour calculer des termes, simuler un phénomène ou déterminer un rang.

\begin{examplebox}{Algorithme de calcul}{algo-simple}
La suite $(u_n)$ est définie par $u_0=4$ et $u_{n+1}=2u_n+3$.

Le programme suivant affiche les termes de $u_0$ à $u_{10}$ :
\begin{verbatim}
u = 4
for n in range(11):
    print(n, u)
    u = 2*u + 3
\end{verbatim}
L'ordre des instructions est important : on affiche d'abord le terme actuel, puis on calcule le suivant.
\end{examplebox}

\begin{exercisebox}{Calculs élémentaires}{exo-1}
Pour chacune des suites suivantes, calculer les quatre premiers termes.
\begin{enumerate}[label=\alph*)]
  \item $u_n=5n-2$ pour $n\in\N$.
  \item $v_n=n^2+n+1$ pour $n\in\N$.
  \item $w_0=3$ et $w_{n+1}=w_n+7$.
  \item $a_0=10$ et $a_{n+1}=2a_n-4$.
\end{enumerate}
\end{exercisebox}

\begin{exercisebox}{Interprétation d'une suite}{exo-2}
Un abonnement coûte $20$ euros le premier mois, puis augmente de $2$ euros chaque mois. On note $p_n$ le prix payé le $n$-ième mois, avec $p_1=20$.
\begin{enumerate}
  \item Que représentent $p_1$, $p_2$ et $p_3$ ?
  \item Donner une relation de récurrence vérifiée par la suite $(p_n)$.
  \item Calculer $p_6$.
  \item Proposer une formule explicite pour $p_n$.
\end{enumerate}
\end{exercisebox}

\section{Sens de variation d'une suite}

\subsection{Suite croissante, décroissante, constante}

\begin{definitionbox}{Sens de variation}{variation}
Soit $(u_n)$ une suite définie sur $\N$.
\begin{itemize}
  \item La suite $(u_n)$ est \textbf{croissante} si, pour tout $n\in\N$,
  \[
  u_{n+1}\geq u_n.
  \]
  \item La suite $(u_n)$ est \textbf{strictement croissante} si, pour tout $n\in\N$,
  \[
  u_{n+1}>u_n.
  \]
  \item La suite $(u_n)$ est \textbf{décroissante} si, pour tout $n\in\N$,
  \[
  u_{n+1}\leq u_n.
  \]
  \item La suite $(u_n)$ est \textbf{strictement décroissante} si, pour tout $n\in\N$,
  \[
  u_{n+1}<u_n.
  \]
  \item La suite $(u_n)$ est \textbf{constante} si, pour tout $n\in\N$,
  \[
  u_{n+1}=u_n.
  \]
\end{itemize}
\end{definitionbox}

\begin{rememberbox}
Étudier les variations d'une suite revient souvent à étudier le signe de la différence :
\[
 u_{n+1}-u_n.
\]
Si cette différence est toujours positive, la suite est croissante. Si elle est toujours négative, la suite est décroissante.
\end{rememberbox}

\subsection{Méthode par différence}

\begin{methodbox}[Étudier les variations avec $u_{n+1}-u_n$]
Pour étudier le sens de variation d'une suite $(u_n)$ :
\begin{enumerate}
  \item calculer $u_{n+1}$ ;
  \item former la différence $u_{n+1}-u_n$ ;
  \item simplifier l'expression ;
  \item déterminer son signe ;
  \item conclure sur les variations de la suite.
\end{enumerate}
\end{methodbox}

\begin{examplebox}{Variation d'une suite explicite}{variation-explicite}
Soit $(u_n)$ définie par
\[
 u_n=2n^2-3n+1.
\]
Étudions son sens de variation.

On calcule :
\[
 u_{n+1}=2(n+1)^2-3(n+1)+1.
\]
Donc
\[
 u_{n+1}=2(n^2+2n+1)-3n-3+1=2n^2+4n+2-3n-2=2n^2+n.
\]
Or
\[
 u_n=2n^2-3n+1.
\]
Ainsi,
\[
 u_{n+1}-u_n=(2n^2+n)-(2n^2-3n+1)=4n-1.
\]
Pour $n\in\N$, on a :
\[
4n-1\geq -1.
\]
Plus précisément, pour $n=0$, $4n-1=-1<0$, et pour $n\geq1$, $4n-1>0$.

Donc $u_1<u_0$, puis la suite est strictement croissante à partir du rang $1$.
\end{examplebox}

\begin{warningbox}
Il est faux de conclure qu'une suite est croissante parce que ses premiers termes augmentent. Les premiers termes permettent de conjecturer, mais pas de démontrer.
\end{warningbox}

\subsection{Méthode par quotient pour les suites positives}

Lorsque les termes d'une suite sont strictement positifs, on peut parfois étudier les variations à l'aide du quotient $\frac{u_{n+1}}{u_n}$.

\begin{propositionbox}{Variation par quotient}{quotient-variation}
Soit $(u_n)$ une suite telle que $u_n>0$ pour tout $n$.
\begin{itemize}
  \item Si $\dfrac{u_{n+1}}{u_n}\geq1$ pour tout $n$, alors $(u_n)$ est croissante.
  \item Si $\dfrac{u_{n+1}}{u_n}\leq1$ pour tout $n$, alors $(u_n)$ est décroissante.
\end{itemize}
\end{propositionbox}

\begin{proofbox}
Supposons $u_n>0$ et $\frac{u_{n+1}}{u_n}\geq1$.
Comme $u_n>0$, on peut multiplier l'inégalité par $u_n$ sans changer le sens de l'inégalité :
\[
 u_{n+1}\geq u_n.
\]
Donc la suite est croissante. La démonstration du cas décroissant est analogue.
\end{proofbox}

\begin{exercisebox}{Variations}{exo-variation}
Étudier le sens de variation des suites suivantes :
\begin{enumerate}[label=\alph*)]
  \item $u_n=4n+7$.
  \item $v_n=-3n+10$.
  \item $w_n=n^2-6n$.
  \item $a_n=2^n$.
  \item $b_n=\left(\dfrac{1}{3}\right)^n$.
\end{enumerate}
\end{exercisebox}

\section{Suites arithmétiques}

\subsection{Définition}

\begin{definitionbox}{Suite arithmétique}{arithmetique}
Une suite $(u_n)$ est dite \textbf{arithmétique} lorsqu'il existe un réel $r$ tel que, pour tout $n$,
\[
 u_{n+1}=u_n+r.
\]
Le nombre $r$ s'appelle la \textbf{raison} de la suite arithmétique.
\end{definitionbox}

\begin{rememberbox}
Une suite arithmétique modélise une évolution où l'on ajoute toujours la même quantité.
\end{rememberbox}

\begin{examplebox}{Reconnaître une suite arithmétique}{reconnaitre-arith}
La suite
\[
5,\;8,\;11,\;14,\;17,\ldots
\]
est arithmétique de raison $3$, car on passe d'un terme au suivant en ajoutant toujours $3$.

La suite
\[
10,\;7,\;4,\;1,\;-2,\ldots
\]
est arithmétique de raison $-3$.
\end{examplebox}

\subsection{Formule explicite}

\begin{theorembox}{Formule explicite d'une suite arithmétique}{formule-arith}
Soit $(u_n)$ une suite arithmétique de raison $r$.
Si elle est définie à partir du rang $0$, alors pour tout $n\in\N$,
\[
 u_n=u_0+nr.
\]
Si elle est définie à partir du rang $1$, alors pour tout $n\geq1$,
\[
 u_n=u_1+(n-1)r.
\]
Plus généralement, pour tous entiers $n$ et $p$ de l'ensemble de définition,
\[
 u_n=u_p+(n-p)r.
\]
\end{theorembox}

\begin{proofbox}[Cas d'une suite commençant au rang $0$]
Par définition d'une suite arithmétique,
\[
 u_{n+1}=u_n+r.
\]
Donc :
\[
 u_1=u_0+r,
\]
\[
 u_2=u_1+r=(u_0+r)+r=u_0+2r,
\]
\[
 u_3=u_2+r=(u_0+2r)+r=u_0+3r.
\]
On observe qu'après $n$ passages, on a ajouté $r$ exactement $n$ fois. Ainsi :
\[
 u_n=u_0+nr.
\]
Une démonstration plus formelle peut se faire par récurrence, mais l'idée essentielle est que l'on ajoute toujours la même quantité.
\end{proofbox}

\begin{examplebox}{Calculer un terme d'une suite arithmétique}{terme-arith}
Soit $(u_n)$ une suite arithmétique telle que $u_0=12$ et de raison $r=5$.
Alors :
\[
 u_n=12+5n.
\]
Donc :
\[
 u_{20}=12+5\times20=112.
\]
\end{examplebox}

\begin{examplebox}{Attention au rang initial}{rang-arith}
Soit $(v_n)$ une suite arithmétique telle que $v_1=12$ et de raison $5$.
Alors :
\[
 v_n=v_1+(n-1)r=12+5(n-1).
\]
Donc :
\[
 v_{20}=12+5\times19=107.
\]
Ce résultat est différent du précédent, car la suite ne commence pas au même rang.
\end{examplebox}

\begin{methodbox}[Montrer qu'une suite est arithmétique]
Pour prouver qu'une suite $(u_n)$ est arithmétique :
\begin{enumerate}
  \item calculer $u_{n+1}-u_n$ ;
  \item montrer que cette différence est constante, c'est-à-dire indépendante de $n$ ;
  \item conclure : la suite est arithmétique et sa raison est cette constante.
\end{enumerate}
\end{methodbox}

\begin{examplebox}{Démontrer qu'une suite est arithmétique}{demo-arith}
Soit $(u_n)$ définie par $u_n=7n-4$.
Montrons que $(u_n)$ est arithmétique.

On calcule :
\[
 u_{n+1}=7(n+1)-4=7n+3.
\]
Donc :
\[
 u_{n+1}-u_n=(7n+3)-(7n-4)=7.
\]
La différence $u_{n+1}-u_n$ est constante. Donc $(u_n)$ est arithmétique de raison $7$.
\end{examplebox}

\subsection{Sens de variation d'une suite arithmétique}

\begin{propositionbox}{Variation d'une suite arithmétique}{variation-arith}
Soit $(u_n)$ une suite arithmétique de raison $r$.
\begin{itemize}
  \item Si $r>0$, alors $(u_n)$ est strictement croissante.
  \item Si $r=0$, alors $(u_n)$ est constante.
  \item Si $r<0$, alors $(u_n)$ est strictement décroissante.
\end{itemize}
\end{propositionbox}

\begin{proofbox}
Par définition, pour tout $n$,
\[
 u_{n+1}-u_n=r.
\]
Si $r>0$, alors $u_{n+1}-u_n>0$, donc $u_{n+1}>u_n$ : la suite est strictement croissante.

Si $r=0$, alors $u_{n+1}=u_n$ : la suite est constante.

Si $r<0$, alors $u_{n+1}-u_n<0$, donc $u_{n+1}<u_n$ : la suite est strictement décroissante.
\end{proofbox}

\subsection{Somme des termes d'une suite arithmétique}

Même si le programme insiste surtout sur la somme des termes d'une suite géométrique, il est très utile de connaître la somme des termes consécutifs d'une suite arithmétique, car elle donne du sens aux calculs et prépare aux raisonnements ultérieurs.

\begin{propositionbox}{Somme de termes consécutifs d'une suite arithmétique}{somme-arith}
La somme de termes consécutifs d'une suite arithmétique est égale à :
\[
\text{nombre de termes}\times \frac{\text{premier terme}+\text{dernier terme}}{2}.
\]
En particulier,
\[
 u_0+u_1+\cdots+u_n=(n+1)\frac{u_0+u_n}{2}.
\]
\end{propositionbox}

\begin{proofbox}[Idée de la démonstration]
Notons
\[
S=u_0+u_1+\cdots+u_n.
\]
Écrivons cette somme dans l'autre sens :
\[
S=u_n+u_{n-1}+\cdots+u_0.
\]
En additionnant ligne à ligne, on obtient :
\[
2S=(u_0+u_n)+(u_1+u_{n-1})+\cdots+(u_n+u_0).
\]
Dans une suite arithmétique, chaque parenthèse vaut $u_0+u_n$. Il y a $n+1$ parenthèses. Donc :
\[
2S=(n+1)(u_0+u_n),
\]
puis :
\[
S=(n+1)\frac{u_0+u_n}{2}.
\]
\end{proofbox}

\begin{exercisebox}{Suites arithmétiques}{exo-arith}
\begin{enumerate}
  \item Soit $(u_n)$ arithmétique de premier terme $u_0=9$ et de raison $-2$.
  \begin{enumerate}
    \item Donner la formule explicite de $u_n$.
    \item Calculer $u_{15}$.
    \item Étudier le sens de variation de $(u_n)$.
  \end{enumerate}
  \item On considère la suite $(v_n)$ définie par $v_n=4n+11$.
  \begin{enumerate}
    \item Montrer que $(v_n)$ est arithmétique.
    \item Donner sa raison et son premier terme.
  \end{enumerate}
  \item Une entreprise produit $500$ pièces le premier jour, puis augmente sa production de $35$ pièces chaque jour.
  \begin{enumerate}
    \item Modéliser la production du jour $n$ par une suite arithmétique.
    \item Calculer la production du 20e jour.
    \item Calculer le nombre total de pièces produites pendant les 20 premiers jours.
  \end{enumerate}
\end{enumerate}
\end{exercisebox}

\section{Suites géométriques}

\subsection{Définition}

\begin{definitionbox}{Suite géométrique}{geometrique}
Une suite $(u_n)$ est dite \textbf{géométrique} lorsqu'il existe un réel $q$ tel que, pour tout $n$,
\[
 u_{n+1}=q u_n.
\]
Le nombre $q$ s'appelle la \textbf{raison} de la suite géométrique.
\end{definitionbox}

\begin{rememberbox}
Une suite géométrique modélise une évolution où l'on multiplie toujours par le même nombre.
C'est le modèle naturel des évolutions en pourcentage constant.
\end{rememberbox}

\begin{examplebox}{Reconnaître une suite géométrique}{reconnaitre-geo}
La suite
\[
3,\;6,\;12,\;24,\;48,\ldots
\]
est géométrique de raison $2$.

La suite
\[
100,\;80,\;64,\;51{,}2,\ldots
\]
est géométrique de raison $0{,}8$.
Elle correspond à une baisse de $20\%$ à chaque étape.
\end{examplebox}

\subsection{Formule explicite}

\begin{theorembox}{Formule explicite d'une suite géométrique}{formule-geo}
Soit $(u_n)$ une suite géométrique de raison $q$.
Si elle est définie à partir du rang $0$, alors pour tout $n\in\N$,
\[
 u_n=u_0 q^n.
\]
Si elle est définie à partir du rang $1$, alors pour tout $n\geq1$,
\[
 u_n=u_1 q^{n-1}.
\]
Plus généralement, pour tous entiers $n$ et $p$ de l'ensemble de définition,
\[
 u_n=u_p q^{n-p}.
\]
\end{theorembox}

\begin{proofbox}[Cas d'une suite commençant au rang $0$]
Par définition d'une suite géométrique,
\[
 u_{n+1}=q u_n.
\]
Donc :
\[
 u_1=q u_0,
\]
\[
 u_2=q u_1=q(q u_0)=q^2u_0,
\]
\[
 u_3=q u_2=q(q^2u_0)=q^3u_0.
\]
Après $n$ passages, on a multiplié $u_0$ par $q$ exactement $n$ fois. Ainsi :
\[
 u_n=u_0q^n.
\]
\end{proofbox}

\begin{methodbox}[Montrer qu'une suite est géométrique]
Pour prouver qu'une suite $(u_n)$ est géométrique :
\begin{enumerate}
  \item vérifier que les termes considérés ne posent pas de problème de division si l'on utilise un quotient ;
  \item calculer $\dfrac{u_{n+1}}{u_n}$ lorsque $u_n\neq0$ ;
  \item montrer que ce quotient est constant ;
  \item conclure : la suite est géométrique et sa raison est cette constante.
\end{enumerate}
On peut aussi montrer directement qu'il existe un réel $q$ tel que $u_{n+1}=q u_n$.
\end{methodbox}

\begin{examplebox}{Démontrer qu'une suite est géométrique}{demo-geo}
Soit $(u_n)$ définie par $u_n=5\times 3^n$.
On calcule :
\[
 u_{n+1}=5\times 3^{n+1}=5\times 3^n\times 3=3u_n.
\]
Donc $(u_n)$ est géométrique de raison $3$.
\end{examplebox}

\begin{warningbox}
Pour une suite géométrique, on ne doit pas additionner la raison : on multiplie par la raison.
Ainsi, si $u_0=10$ et $q=1{,}05$, alors :
\[
 u_1=10\times1{,}05=10{,}5,
\]
et non $11{,}05$.
\end{warningbox}

\subsection{Pourcentages et coefficient multiplicateur}

Les suites géométriques sont très liées aux pourcentages.

\begin{propositionbox}{Pourcentage d'évolution}{pourcentage}
Augmenter une quantité de $t\%$ revient à la multiplier par :
\[
1+\frac{t}{100}.
\]
Diminuer une quantité de $t\%$ revient à la multiplier par :
\[
1-\frac{t}{100}.
\]
\end{propositionbox}

\begin{examplebox}{Placement financier}{placement-simple}
On place $1\,000$ euros sur un compte rémunéré à $3\%$ par an. On note $C_n$ le capital au bout de $n$ années.

Chaque année, le capital est multiplié par $1{,}03$. Donc :
\[
 C_{n+1}=1{,}03 C_n.
\]
La suite $(C_n)$ est géométrique de premier terme $C_0=1000$ et de raison $1{,}03$.
Ainsi :
\[
 C_n=1000\times1{,}03^n.
\]
Au bout de $10$ ans :
\[
 C_{10}=1000\times1{,}03^{10}\approx1343{,}92.
\]
Le capital est donc d'environ $1\,343{,}92$ euros.
\end{examplebox}

\subsection{Sens de variation d'une suite géométrique positive}

\begin{propositionbox}{Variation d'une suite géométrique positive}{variation-geo}
Soit $(u_n)$ une suite géométrique de raison $q$ et de premier terme strictement positif.
\begin{itemize}
  \item Si $q>1$, alors $(u_n)$ est strictement croissante.
  \item Si $q=1$, alors $(u_n)$ est constante.
  \item Si $0<q<1$, alors $(u_n)$ est strictement décroissante.
\end{itemize}
\end{propositionbox}

\begin{proofbox}
Comme $u_n>0$, on a :
\[
 u_{n+1}-u_n=qu_n-u_n=(q-1)u_n.
\]
Le signe de $u_{n+1}-u_n$ est donc le signe de $q-1$.
Si $q>1$, la différence est positive : la suite est croissante.
Si $q=1$, la différence est nulle : la suite est constante.
Si $0<q<1$, la différence est négative : la suite est décroissante.
\end{proofbox}

\begin{warningbox}
La proposition précédente suppose que le premier terme est positif et que la raison est positive. Si la raison est négative, les signes alternent et la suite n'est en général ni croissante ni décroissante.
\end{warningbox}

\subsection{Somme des termes d'une suite géométrique}

\begin{theorembox}{Somme des puissances successives}{somme-puissances}
Pour tout réel $q\neq1$ et tout entier naturel $n$,
\[
1+q+q^2+\cdots+q^n=\frac{1-q^{n+1}}{1-q}.
\]
Si $q=1$, alors :
\[
1+1+\cdots+1=n+1.
\]
\end{theorembox}

\begin{proofbox}[Démonstration fondamentale]
Posons
\[
S=1+q+q^2+\cdots+q^n.
\]
Multiplions par $q$ :
\[
qS=q+q^2+q^3+\cdots+q^{n+1}.
\]
Soustrayons :
\[
S-qS=(1+q+q^2+\cdots+q^n)-(q+q^2+\cdots+q^n+q^{n+1}).
\]
Tous les termes intermédiaires se simplifient, donc :
\[
(1-q)S=1-q^{n+1}.
\]
Si $q\neq1$, on divise par $1-q$ :
\[
S=\frac{1-q^{n+1}}{1-q}.
\]
\end{proofbox}

\begin{propositionbox}{Somme de termes consécutifs d'une suite géométrique}{somme-geo}
Soit $(u_n)$ une suite géométrique de raison $q\neq1$.
Alors :
\[
 u_0+u_1+\cdots+u_n=u_0\frac{1-q^{n+1}}{1-q}.
\]
Plus généralement, la somme de $N$ termes consécutifs d'une suite géométrique est :
\[
\text{premier terme}\times\frac{1-q^{N}}{1-q}.
\]
\end{propositionbox}

\begin{examplebox}{Calcul d'une somme géométrique}{somme-geo-exemple}
Calculons :
\[
S=3+6+12+24+48.
\]
C'est la somme de $5$ termes d'une suite géométrique de premier terme $3$ et de raison $2$.
Donc :
\[
S=3\times\frac{1-2^5}{1-2}=3\times\frac{1-32}{-1}=3\times31=93.
\]
\end{examplebox}

\begin{warningbox}
Dans la formule
\[
\text{premier terme}\times\frac{1-q^N}{1-q},
\]
le nombre $N$ désigne le nombre de termes, et non le dernier rang.
Si l'on somme de $u_0$ à $u_n$, il y a $n+1$ termes.
\end{warningbox}

\begin{exercisebox}{Suites géométriques}{exo-geo}
\begin{enumerate}
  \item Soit $(u_n)$ géométrique de premier terme $u_0=7$ et de raison $2$.
  \begin{enumerate}
    \item Donner la formule explicite de $u_n$.
    \item Calculer $u_{10}$.
    \item Calculer $u_0+u_1+\cdots+u_{10}$.
  \end{enumerate}
  \item Soit $(v_n)$ définie par $v_n=12\times0{,}8^n$.
  \begin{enumerate}
    \item Montrer que $(v_n)$ est géométrique.
    \item Donner sa raison.
    \item Étudier son sens de variation.
  \end{enumerate}
  \item Un objet perd $15\%$ de sa valeur chaque année. Il vaut initialement $800$ euros.
  \begin{enumerate}
    \item Modéliser la situation par une suite.
    \item Donner une formule explicite.
    \item Calculer sa valeur après $6$ ans.
  \end{enumerate}
\end{enumerate}
\end{exercisebox}

\section{Modélisation de phénomènes discrets}

\subsection{Choisir entre modèle arithmétique et modèle géométrique}

\begin{methodbox}[Identifier le bon modèle]
Pour modéliser une situation discrète :
\begin{enumerate}
  \item Si l'on ajoute ou retire toujours la même quantité, on utilise une suite arithmétique.
  \item Si l'on multiplie toujours par le même coefficient, notamment lors d'une évolution en pourcentage constant, on utilise une suite géométrique.
  \item On précise toujours le rang initial et la signification de $u_n$.
\end{enumerate}
\end{methodbox}

\begin{center}
\begin{tabularx}{0.95\textwidth}{|>{\bfseries}p{4cm}|X|X|}
\hline
Situation & Type de suite & Exemple \\
\hline
On ajoute $50$ chaque année & Arithmétique & $u_{n+1}=u_n+50$ \\
\hline
On retire $12$ chaque mois & Arithmétique & $u_{n+1}=u_n-12$ \\
\hline
On augmente de $4\%$ par an & Géométrique & $u_{n+1}=1{,}04u_n$ \\
\hline
On baisse de $7\%$ par semaine & Géométrique & $u_{n+1}=0{,}93u_n$ \\
\hline
\end{tabularx}
\end{center}

\subsection{Évolution d'une population}

\begin{examplebox}{Population à augmentation constante}{population-arith}
Une ville compte $25\,000$ habitants en 2025. Chaque année, elle gagne environ $600$ habitants. On note $P_n$ la population $n$ années après 2025.

On a :
\[
P_0=25\,000,
\qquad
P_{n+1}=P_n+600.
\]
La suite est arithmétique de raison $600$.
Donc :
\[
P_n=25\,000+600n.
\]
En 2035, soit $10$ ans après 2025 :
\[
P_{10}=25\,000+600\times10=31\,000.
\]
\end{examplebox}

\begin{examplebox}{Population à taux constant}{population-geo}
Une population de bactéries augmente de $12\%$ toutes les heures. Elle compte initialement $5\,000$ bactéries. On note $B_n$ le nombre de bactéries après $n$ heures.

Une augmentation de $12\%$ correspond à un coefficient multiplicateur :
\[
1+\frac{12}{100}=1{,}12.
\]
Donc :
\[
B_{n+1}=1{,}12B_n.
\]
La suite est géométrique de raison $1{,}12$ et :
\[
B_n=5000\times1{,}12^n.
\]
Après $8$ heures :
\[
B_8=5000\times1{,}12^8\approx12\,380.
\]
\end{examplebox}

\subsection{Placement financier}

\begin{examplebox}{Capital avec intérêts composés}{interets-composes}
On place $2\,500$ euros à intérêts composés au taux annuel de $4\%$. On note $C_n$ le capital après $n$ années.

Chaque année :
\[
C_{n+1}=1{,}04C_n.
\]
Donc :
\[
C_n=2500\times1{,}04^n.
\]
Après $15$ ans :
\[
C_{15}=2500\times1{,}04^{15}\approx4502{,}36.
\]
\end{examplebox}

\begin{examplebox}{Capital avec versement constant}{capital-versement}
On dépose $100$ euros chaque mois sur un compte, sans intérêt. Le capital augmente donc de $100$ euros par mois. Si $C_0=500$, alors :
\[
C_{n+1}=C_n+100.
\]
La suite est arithmétique et :
\[
C_n=500+100n.
\]
\end{examplebox}

\subsection{Phénomènes répétés}

\begin{examplebox}{Rebond d'une balle}{rebond}
Une balle rebondit à $70\%$ de la hauteur précédente. Elle est lâchée d'une hauteur initiale de $2$ mètres. On note $h_n$ la hauteur atteinte après le $n$-ième rebond, avec $h_0=2$.

On a :
\[
h_{n+1}=0{,}7h_n.
\]
Donc :
\[
h_n=2\times0{,}7^n.
\]
Après $5$ rebonds :
\[
h_5=2\times0{,}7^5\approx0{,}336.
\]
La hauteur est donc d'environ $34$ cm.
\end{examplebox}

\begin{exercisebox}{Modélisation}{exo-modelisation}
\begin{enumerate}
  \item Une population de poissons dans un lac diminue de $8\%$ par an. Elle est initialement de $12\,000$ poissons.
  \begin{enumerate}
    \item Définir une suite adaptée.
    \item Donner la relation de récurrence.
    \item Donner la formule explicite.
    \item Estimer la population après $10$ ans.
  \end{enumerate}
  \item Un club compte $80$ adhérents et gagne $15$ adhérents par mois.
  \begin{enumerate}
    \item Modéliser par une suite.
    \item Calculer le nombre d'adhérents après $18$ mois.
  \end{enumerate}
  \item Un salaire mensuel de $1\,800$ euros augmente de $2\%$ chaque année.
  \begin{enumerate}
    \item Modéliser la situation.
    \item Calculer le salaire après $5$ ans.
    \item Déterminer à partir de quelle année le salaire dépasse $2\,000$ euros.
  \end{enumerate}
\end{enumerate}
\end{exercisebox}

\section{Algorithmique et suites}

\subsection{Calculer un terme}

\begin{methodbox}[Calculer $u_n$ avec Python]
Pour calculer un terme d'une suite définie par récurrence :
\begin{enumerate}
  \item initialiser la variable avec le premier terme ;
  \item répéter la relation de récurrence le bon nombre de fois ;
  \item afficher ou retourner le résultat.
\end{enumerate}
\end{methodbox}

\begin{examplebox}{Programme Python}{python-terme}
Pour $u_0=2$ et $u_{n+1}=1{,}5u_n+4$, on peut calculer $u_{20}$ avec :
\begin{verbatim}
u = 2
for k in range(20):
    u = 1.5*u + 4
print(u)
\end{verbatim}
Après la boucle, la variable contient $u_{20}$.
\end{examplebox}

\subsection{Déterminer un seuil}

\begin{examplebox}{Recherche d'un seuil}{python-seuil}
On place $1\,000$ euros à $5\%$ par an. On cherche au bout de combien d'années le capital dépasse $2\,000$ euros.

On peut écrire :
\begin{verbatim}
C = 1000
n = 0
while C <= 2000:
    C = 1.05*C
    n = n + 1
print(n)
\end{verbatim}
La variable $n$ donne le premier rang pour lequel $C_n>2000$.
\end{examplebox}

\begin{warningbox}
Dans une boucle \texttt{while}, il faut vérifier que la condition finira par devenir fausse. Sinon, la boucle ne s'arrête jamais.
\end{warningbox}

\section{Exercices bilan de chapitre}

\begin{exercisebox}{Bilan 1 -- Nature d'une suite}{bilan-1}
Pour chacune des suites suivantes, dire si elle est arithmétique, géométrique ou ni l'une ni l'autre. Justifier.
\begin{enumerate}
  \item $u_n=3n+8$.
  \item $v_n=5\times2^n$.
  \item $w_n=n^2+1$.
  \item $a_0=4$ et $a_{n+1}=a_n-6$.
  \item $b_0=7$ et $b_{n+1}=0{,}9b_n$.
\end{enumerate}
\end{exercisebox}

\begin{exercisebox}{Bilan 2 -- Sommes}{bilan-2}
Calculer les sommes suivantes :
\begin{enumerate}
  \item $1+2+4+8+\cdots+2^{10}$.
  \item $5+15+45+\cdots+5\times3^8$.
  \item $100+90+81+\cdots+100\times0{,}9^{12}$.
\end{enumerate}
\end{exercisebox}

\begin{exercisebox}{Bilan 3 -- Problème de modélisation}{bilan-3}
Une entreprise vend $12\,000$ unités d'un produit en 2026. Elle prévoit une augmentation annuelle de $6\%$ pendant plusieurs années.
\begin{enumerate}
  \item Définir une suite modélisant les ventes.
  \item Donner sa nature et sa raison.
  \item Donner une formule explicite.
  \item Estimer les ventes en 2031.
  \item Calculer le nombre total d'unités vendues de 2026 à 2031 inclus.
\end{enumerate}
\end{exercisebox}

\begin{exercisebox}{Bilan 4 -- Comparaison de modèles}{bilan-4}
Deux placements sont proposés.
\begin{itemize}
  \item Placement A : capital initial $2\,000$ euros, gain fixe de $120$ euros par an.
  \item Placement B : capital initial $2\,000$ euros, taux annuel de $5\%$.
\end{itemize}
On note $A_n$ et $B_n$ les capitaux après $n$ années.
\begin{enumerate}
  \item Donner les formules explicites de $A_n$ et $B_n$.
  \item Calculer les capitaux après $5$ ans, puis après $10$ ans.
  \item À l'aide d'un raisonnement ou d'un algorithme, déterminer à partir de quelle année le placement B devient plus avantageux que le placement A.
\end{enumerate}
\end{exercisebox}

\newpage
\section{Grand devoir bilan final}

\begin{center}
\Large\bfseries Devoir bilan -- Suites numériques\end{center}

\begin{exercisebox}{Exercice 1 -- Questions de cours}{ds-1}
\begin{enumerate}
  \item Donner la définition d'une suite arithmétique.
  \item Donner la formule explicite d'une suite arithmétique de premier terme $u_0$ et de raison $r$.
  \item Donner la définition d'une suite géométrique.
  \item Donner la formule de la somme $1+q+q^2+\cdots+q^n$ lorsque $q\neq1$.
\end{enumerate}
\end{exercisebox}

\begin{exercisebox}{Exercice 2 -- Étude d'une suite arithmétique}{ds-2}
On considère la suite $(u_n)$ définie par $u_0=18$ et, pour tout $n\in\N$,
\[
 u_{n+1}=u_n-4.
\]
\begin{enumerate}
  \item Calculer $u_1$, $u_2$ et $u_3$.
  \item Justifier que $(u_n)$ est arithmétique et préciser sa raison.
  \item Donner une formule explicite de $u_n$.
  \item Calculer $u_{25}$.
  \item Déterminer le plus petit entier $n$ tel que $u_n<0$.
\end{enumerate}
\end{exercisebox}

\begin{exercisebox}{Exercice 3 -- Étude d'une suite géométrique}{ds-3}
On considère la suite $(v_n)$ définie par $v_0=500$ et, pour tout $n\in\N$,
\[
 v_{n+1}=1{,}08v_n.
\]
\begin{enumerate}
  \item Donner la nature de la suite et sa raison.
  \item Donner une formule explicite de $v_n$.
  \item Calculer $v_6$ à l'euro près.
  \item Calculer $v_0+v_1+\cdots+v_6$ à l'euro près.
\end{enumerate}
\end{exercisebox}

\begin{exercisebox}{Exercice 4 -- Modélisation}{ds-4}
Une association compte $300$ membres au 1er janvier 2026. Chaque année, elle perd $5\%$ de ses membres, mais gagne ensuite $20$ nouveaux membres.
On note $m_n$ le nombre de membres au bout de $n$ années, avec $m_0=300$.
\begin{enumerate}
  \item Expliquer pourquoi la relation de récurrence est
  \[
  m_{n+1}=0{,}95m_n+20.
  \]
  \item Calculer $m_1$, $m_2$ et $m_3$.
  \item Cette suite est-elle arithmétique ? géométrique ? Justifier.
  \item On pose $a_n=m_n-400$. Montrer que $(a_n)$ est géométrique.
  \item En déduire une formule explicite de $m_n$.
\end{enumerate}
\end{exercisebox}

\newpage
\section{Corrections détaillées}

\subsection*{Correction des exercices de cours}
\addcontentsline{toc}{section}{Corrections détaillées}

\begin{correctionbox}[Exercice \ref{exo:exo-1}]
\begin{enumerate}[label=\alph*)]
  \item $u_0=-2$, $u_1=3$, $u_2=8$, $u_3=13$.
  \item $v_0=1$, $v_1=3$, $v_2=7$, $v_3=13$.
  \item $w_0=3$, $w_1=10$, $w_2=17$, $w_3=24$.
  \item $a_0=10$, $a_1=16$, $a_2=28$, $a_3=52$.
\end{enumerate}
\end{correctionbox}

\begin{correctionbox}[Exercice \ref{exo:exo-2}]
\begin{enumerate}
  \item $p_1$ est le prix du premier mois, $p_2$ celui du deuxième mois, $p_3$ celui du troisième mois.
  \item Comme le prix augmente de $2$ euros chaque mois : $p_{n+1}=p_n+2$.
  \item $p_6=20+5\times2=30$ euros.
  \item Comme la suite commence au rang $1$, $p_n=20+2(n-1)$.
\end{enumerate}
\end{correctionbox}

\begin{correctionbox}[Exercice \ref{exo:exo-variation}]
\begin{enumerate}[label=\alph*)]
  \item $u_{n+1}-u_n=4>0$, donc $(u_n)$ est strictement croissante.
  \item $v_{n+1}-v_n=-3<0$, donc $(v_n)$ est strictement décroissante.
  \item $w_{n+1}-w_n=(n+1)^2-6(n+1)-(n^2-6n)=2n-5$. La suite décroît jusqu'au voisinage du rang $2$, puis croît à partir du rang $3$.
  \item $a_n=2^n$ est géométrique de raison $2>1$ et de premier terme positif, donc croissante.
  \item $b_n=(1/3)^n$ est géométrique de raison $1/3$, comprise entre $0$ et $1$, et de premier terme positif, donc décroissante.
\end{enumerate}
\end{correctionbox}

\begin{correctionbox}[Exercice \ref{exo:exo-arith}]
\begin{enumerate}
  \item $u_n=9-2n$. Ainsi $u_{15}=9-30=-21$. Comme la raison est négative, la suite est strictement décroissante.
  \item $v_{n+1}-v_n=[4(n+1)+11]-(4n+11)=4$. La suite est arithmétique de raison $4$ et $v_0=11$.
  \item Si $p_n$ désigne la production du $n$-ième jour avec $p_1=500$, alors $p_n=500+35(n-1)$. Le 20e jour : $p_{20}=500+35\times19=1165$. Le total sur 20 jours vaut $20\times\frac{500+1165}{2}=16650$ pièces.
\end{enumerate}
\end{correctionbox}

\begin{correctionbox}[Exercice \ref{exo:exo-geo}]
\begin{enumerate}
  \item $u_n=7\times2^n$. Donc $u_{10}=7\times1024=7168$. La somme de $u_0$ à $u_{10}$ vaut $7\frac{1-2^{11}}{1-2}=7(2047)=14329$.
  \item $v_{n+1}=12\times0{,}8^{n+1}=0{,}8v_n$. La suite est géométrique de raison $0{,}8$. Elle est décroissante car $0<0{,}8<1$ et $v_0>0$.
  \item Une baisse de $15\%$ correspond à un coefficient $0{,}85$. Si $V_0=800$, alors $V_n=800\times0{,}85^n$. Après $6$ ans : $V_6\approx301{,}70$ euros.
\end{enumerate}
\end{correctionbox}

\begin{correctionbox}[Exercice \ref{exo:exo-modelisation}]
\begin{enumerate}
  \item Une baisse de $8\%$ correspond à un coefficient $0{,}92$. Donc $P_0=12000$ et $P_{n+1}=0{,}92P_n$, d'où $P_n=12000\times0{,}92^n$. Après $10$ ans : $P_{10}\approx5218$ poissons.
  \item $A_0=80$ et $A_{n+1}=A_n+15$, donc $A_n=80+15n$. Après $18$ mois : $A_{18}=350$.
  \item $S_0=1800$ et $S_{n+1}=1{,}02S_n$, donc $S_n=1800\times1{,}02^n$. Après $5$ ans : $S_5\approx1987{,}35$. Pour dépasser $2000$, on teste : $S_5<2000$ et $S_6\approx2027{,}10$, donc à partir de la 6e année.
\end{enumerate}
\end{correctionbox}

\subsection*{Correction des exercices bilan}

\begin{correctionbox}[Bilan 1]
\begin{enumerate}
  \item $u_n=3n+8$ est arithmétique de raison $3$.
  \item $v_n=5\times2^n$ est géométrique de raison $2$.
  \item $w_n=n^2+1$ n'est ni arithmétique ni géométrique.
  \item $a_{n+1}=a_n-6$ : arithmétique de raison $-6$.
  \item $b_{n+1}=0{,}9b_n$ : géométrique de raison $0{,}9$.
\end{enumerate}
\end{correctionbox}

\begin{correctionbox}[Bilan 2]
\begin{enumerate}
  \item $1+2+\cdots+2^{10}=\frac{1-2^{11}}{1-2}=2047$.
  \item $5+15+\cdots+5\times3^8=5\frac{1-3^9}{1-3}=49205$.
  \item $100+90+\cdots+100\times0{,}9^{12}=100\frac{1-0{,}9^{13}}{1-0{,}9}\approx745{,}81$.
\end{enumerate}
\end{correctionbox}

\begin{correctionbox}[Bilan 3]
On note $V_n$ les ventes $n$ années après 2026. Alors $V_0=12000$ et $V_{n+1}=1{,}06V_n$. La suite est géométrique de raison $1{,}06$ et $V_n=12000\times1{,}06^n$.

L'année 2031 correspond à $n=5$, donc $V_5=12000\times1{,}06^5\approx16059$.

Le total de 2026 à 2031 inclus correspond à $V_0+\cdots+V_5$ :
\[
12000\frac{1-1{,}06^6}{1-1{,}06}\approx83704.
\]
\end{correctionbox}

\begin{correctionbox}[Bilan 4]
On a :
\[
A_n=2000+120n,
\qquad
B_n=2000\times1{,}05^n.
\]
Après $5$ ans : $A_5=2600$ et $B_5\approx2552{,}56$. Après $10$ ans : $A_{10}=3200$ et $B_{10}\approx3257{,}79$.

On teste les rangs ou on écrit un programme. On trouve que $B_n>A_n$ pour la première fois à partir de $n=9$ :
\[
B_8\approx2954{,}91<2960=A_8,
\]
mais
\[
B_9\approx3102{,}66>3080=A_9.
\]
\end{correctionbox}

\subsection*{Correction du devoir bilan final}

\begin{correctionbox}[Exercice 1]
\begin{enumerate}
  \item Une suite est arithmétique s'il existe un réel $r$ tel que $u_{n+1}=u_n+r$ pour tout $n$.
  \item Si la suite commence au rang $0$, $u_n=u_0+nr$.
  \item Une suite est géométrique s'il existe un réel $q$ tel que $u_{n+1}=qu_n$ pour tout $n$.
  \item Pour $q\neq1$, $1+q+\cdots+q^n=\frac{1-q^{n+1}}{1-q}$.
\end{enumerate}
\end{correctionbox}

\begin{correctionbox}[Exercice 2]
\begin{enumerate}
  \item $u_1=14$, $u_2=10$, $u_3=6$.
  \item La relation $u_{n+1}=u_n-4$ montre que la suite est arithmétique de raison $-4$.
  \item $u_n=18-4n$.
  \item $u_{25}=18-100=-82$.
  \item On cherche $18-4n<0$, soit $18<4n$, donc $n>4{,}5$. Le plus petit entier est $n=5$.
\end{enumerate}
\end{correctionbox}

\begin{correctionbox}[Exercice 3]
\begin{enumerate}
  \item La suite est géométrique de raison $1{,}08$.
  \item $v_n=500\times1{,}08^n$.
  \item $v_6=500\times1{,}08^6\approx793$ euros.
  \item $v_0+\cdots+v_6=500\frac{1-1{,}08^7}{1-1{,}08}\approx4461$ euros.
\end{enumerate}
\end{correctionbox}

\begin{correctionbox}[Exercice 4]
\begin{enumerate}
  \item Perdre $5\%$ revient à multiplier par $0{,}95$, puis on ajoute $20$ membres. Donc $m_{n+1}=0{,}95m_n+20$.
  \item $m_1=0{,}95\times300+20=305$. Puis $m_2=0{,}95\times305+20=309{,}75$. Enfin $m_3=0{,}95\times309{,}75+20=314{,}2625$.
  \item La suite n'est pas arithmétique car les différences ne sont pas constantes. Elle n'est pas géométrique car les quotients ne sont pas constants.
  \item $a_{n+1}=m_{n+1}-400=0{,}95m_n+20-400=0{,}95m_n-380=0{,}95(m_n-400)=0{,}95a_n$. Donc $(a_n)$ est géométrique de raison $0{,}95$.
  \item $a_0=m_0-400=-100$. Donc $a_n=-100\times0{,}95^n$. Ainsi :
  \[
  m_n=400-100\times0{,}95^n.
  \]
\end{enumerate}
\end{correctionbox}

\newpage
\section{Fiche de synthèse finale}

\begin{rememberbox}[Définitions essentielles]
\begin{itemize}
  \item Une suite est une fonction définie sur des entiers.
  \item Suite explicite : $u_n$ est donné directement en fonction de $n$.
  \item Suite récurrente : $u_{n+1}$ est donné en fonction de $u_n$.
  \item Suite arithmétique : $u_{n+1}=u_n+r$.
  \item Suite géométrique : $u_{n+1}=qu_n$.
\end{itemize}
\end{rememberbox}

\begin{rememberbox}[Formules à connaître]
Pour une suite arithmétique de raison $r$ :
\[
 u_n=u_0+nr.
\]
Pour une suite géométrique de raison $q$ :
\[
 u_n=u_0q^n.
\]
Somme géométrique :
\[
1+q+q^2+\cdots+q^n=\frac{1-q^{n+1}}{1-q}\quad(q\neq1).
\]
Somme de $N$ termes géométriques consécutifs :
\[
\text{premier terme}\times\frac{1-q^N}{1-q}.
\]
\end{rememberbox}

\begin{methodbox}[Réflexes de résolution]
\begin{enumerate}
  \item Identifier le rang initial.
  \item Écrire clairement ce que représente $u_n$.
  \item Distinguer augmentation constante et pourcentage constant.
  \item Pour une suite arithmétique : utiliser les différences.
  \item Pour une suite géométrique : utiliser les quotients ou la relation multiplicative.
  \item Pour une somme géométrique : compter correctement le nombre de termes.
\end{enumerate}
\end{methodbox}

\begin{warningbox}[Erreurs fréquentes]
\begin{itemize}
  \item Confondre $u_n$ et $(u_n)$.
  \item Oublier que de $u_0$ à $u_n$, il y a $n+1$ termes.
  \item Utiliser $u_n=u_0+nr$ alors que la suite commence au rang $1$.
  \item Transformer une augmentation de $5\%$ en ajout de $5$ au lieu d'une multiplication par $1{,}05$.
  \item Conclure sur les variations à partir de quelques termes seulement.
\end{itemize}
\end{warningbox}

\section*{Indications bibliographiques et institutionnelles}
\addcontentsline{toc}{section}{Indications bibliographiques et institutionnelles}
Ce cours suit les attendus usuels de Première générale spécialité mathématiques : suites comme modèles d'évolutions discrètes, calcul de termes définis explicitement ou par récurrence, suites arithmétiques et géométriques, sens de variation et sommes finies de termes géométriques. Il est conseillé au professeur de l'adapter au rythme de la classe, en ajoutant des activités d'introduction et des travaux Python selon les besoins.

\end{document}
